\section{Additional Exercises}

\subsection{\texorpdfstring{Virtuoso Use Of The Symbol $\big(\frac{\p}{\p x^i}\big)_p$}{Virtuoso Use Of The Symbol (d/dx\^i)\_p}}

\begin{exercise}\label{ex:tangentspaces-1}
Show that, for overlapping charts $(U,x)$ and $(V,y)$, one has
\bse
    \bigg(\frac{\p x^a}{\p y^m}\bigg)_p \bigg(\frac{\p y^m}{\p x^b}\bigg)_p = \del^a_b
\ese
for any $p\in U\cap V$.
\end{exercise}
\hyperref[sol:tangentspaces-1]{\small \textit{Solution on p.~\pageref*{sol:tangentspaces-1}}}

\begin{exercise}\label{ex:tangentspaces-2}
After inserting $y^{-1}\circ y$, where $y$ is another chart map on the same chart domain $U$, at the appropriate position in the definition of the left hand side of
\bse
    \bigg(\frac{\p f}{\p x^i}\bigg)_p = \bigg(\frac{\p y^m}{\p x^i}\bigg)_p\bigg(\frac{\p f}{\p y^m}\bigg)_p,
\ese
use the multidimensional chain rule to show that it equals the right-hand side.
\end{exercise}
\hyperref[sol:tangentspaces-2]{\small \textit{Solution on p.~\pageref*{sol:tangentspaces-2}}}

\begin{exercise}\label{ex:tangentspaces-3}
Do the $\dim\cM$ many quantities defined by the left-hand side of the above expression constitute the components of a tensor? If so, what is the valence and rank of the tensor?
\end{exercise}
\hyperref[sol:tangentspaces-3]{\small \textit{Solution on p.~\pageref*{sol:tangentspaces-3}}}

\subsection{Transformation Of Vector Components}

\begin{exercise}\label{ex:tangentspaces-4}
Let the topological space $(\R^2,\cO_{st.})$ be equipped with the atlas $\cA = \{(\R^2,x),(\R^2,y)\}$, where
\bse
    x : (a,b) \mapsto (a,b) \qand y:(a,b) \mapsto (a,b+a^3).
\ese
Calculate the objects $\big(\frac{\p x^i}{\p y^j}\big)_p$!
\end{exercise}
\hyperref[sol:tangentspaces-4]{\small \textit{Solution on p.~\pageref*{sol:tangentspaces-4}}}

\begin{exercise}\label{ex:tangentspaces-5}
(Reworded slightly to save typing) Recall that the components of the velocity to a curve $\gamma$ in a chart $(U,x)$ at point $p=\gamma(\lambda_0)$ are given by
\bse
    \dot{\gamma}_x^i(\lambda_0) := \big(x\circ \gamma)^{i\prime} (\lambda_0).
\ese
Now consider the curve
\bse
    \gamma :\R\to \R^2; \quad  \lambda\mapsto (\lambda,-\lambda).
\ese
Calculate the components $\dot{\gamma}_x^i (\lambda_0)$ and $\dot{\gamma}_y^i (\lambda_0)$!
\end{exercise}
\hyperref[sol:tangentspaces-5]{\small \textit{Solution on p.~\pageref*{sol:tangentspaces-5}}}

\begin{exercise}\label{ex:tangentspaces-6}
With the above results in mind, how could you have obtained the components of $\dot{\gamma}_x^i(\lambda_0)$ from the $\dot{\gamma}_y^i(\lambda_0)$?
\end{exercise}
\hyperref[sol:tangentspaces-6]{\small \textit{Solution on p.~\pageref*{sol:tangentspaces-6}}}

\subsection{The Gradient}

\begin{exercise}\label{ex:tangentspaces-7}
Given a function $f$ on a manifold $\cM$, the level sets of $f$ for a constant $c\in\R$ are defined as
\bse
    N_c(f) := \{p\in\cM \, | \, f(p) = c\}.
\ese
Formulate the condition for a curve $\gamma:\R\to\cM$ to take values solely in one of the level sets of a function $f$!
\end{exercise}
\hyperref[sol:tangentspaces-7]{\small \textit{Solution on p.~\pageref*{sol:tangentspaces-7}}}

\begin{exercise}\label{ex:tangentspaces-8}
Now show that the gradient of the function annihilates the velocity vector $v_{\gamma,p}$ for any such $\gamma$ through $p$ in $N_c(f)$. In other words, show that
\bse
    (df)_p (v_{\gamma,p})=0.
\ese
\end{exercise}
\hyperref[sol:tangentspaces-8]{\small \textit{Solution on p.~\pageref*{sol:tangentspaces-8}}}

\subsection{Is There A Well-Defined Sum Of Curves?}

\begin{exercise}\label{ex:tangentspaces-9}
Let the topological manifold $(\R^2,\cO_{st.})$ be equipped with the atlas $\cA=\{(\R^2,x),(\R^2,y)\}$ where
\bse
    x:(a,b) \mapsto (a,b) \qand y:(a,b)\mapsto (a, b\cdot e^a).
\ese
Is $\cA$ a $C^{\infty}$-atlas?
\end{exercise}
\hyperref[sol:tangentspaces-9]{\small \textit{Solution on p.~\pageref*{sol:tangentspaces-9}}}

\begin{exercise}\label{ex:tangentspaces-10}
On our $\R^2$ above, consider two curves $\gamma,\del:\R\to\R^2$ given by
\bse
    \gamma: \lambda \mapsto (\lambda, 1),  \qand  \del :\lambda \mapsto  (1,\lambda).
\ese
Without referring to any chart, can you give the sum $\gamma + \del$ of these curves?
\end{exercise}
\hyperref[sol:tangentspaces-10]{\small \textit{Solution on p.~\pageref*{sol:tangentspaces-10}}}

\begin{exercise}\label{ex:tangentspaces-11}
Calculate the representatives of both curves with respect to both charts. Illustrate the results. Where do the curves in the charts intersect?
\end{exercise}
\hyperref[sol:tangentspaces-11]{\small \textit{Solution on p.~\pageref*{sol:tangentspaces-11}}}

\begin{exercise}\label{ex:tangentspaces-12}
(Reworded to save typing) Use the formula from the lectures
\bse
    \sig_x: \R\to U; \quad \lambda \mapsto x^{-1}\big( (x\circ\gamma)(\lambda+\lambda_0) + (x\circ\del)(\lambda+\lambda_1) - (x\circ\gamma)(\lambda_0)\big),
\ese
where $\gamma(\lambda_0)=\del(\lambda_1)$, to find the sum $\gamma+\del$. Also do the calculation for $\sig_y(\lambda)$.
\end{exercise}
\hyperref[sol:tangentspaces-12]{\small \textit{Solution on p.~\pageref*{sol:tangentspaces-12}}}

\begin{exercise}\label{ex:tangentspaces-13}
Show that -- desipte the above results -- the velocity of $\sig_x$ and the velocity $\sig_y$ are equal at the intersection point.
\end{exercise}
\hyperref[sol:tangentspaces-13]{\small \textit{Solution on p.~\pageref*{sol:tangentspaces-13}}}
